\section{Introduction}\label{sec:introduction}

Uniform discreteness keeps distinct points of a metric space a fixed positive
distance apart. For its Lipschitz-free space, this separation guarantees the
approximation property: the identity can be approximated on each compact set
by finite-rank operators. Kalton asked whether the approximating operators
can always be chosen with one bound on their norms
\cite[Proposition~4.4 and following remarks, p.~185]{Kalton04}. The question
was restated by Godefroy and Ozawa \cite[Question~1]{GodefroyOzawa14} and in
Godefroy's survey \cite[Problem~6.2]{Godefroy15}. We construct a countable
uniformly discrete space for which no such bound exists.

We work over the real scalars. For a pointed metric space $(M,d,o)$, let
$\Lip_0(M)$ be the Banach space of real Lipschitz functions vanishing at $o$,
with the Lipschitz constant as norm. Its \emph{Lipschitz-free space} is
\[
 \F(M)=\overline{\operatorname{span}}\{\delta_x:x\in M\}
 \subseteq\Lip_0(M)^*,\qquad \delta_x(f)=f(x).
\]
Thus $\delta_o=0$ and $\norm{\delta_x-\delta_y}=d(x,y)$. The latter
identity follows from the Lipschitz inequality and the test function
$z\mapsto d(z,y)-d(o,y)$.
The construction originates in the molecular spaces of Arens and Eells
\cite[\S2]{ArensEells56}; its linearization properties and complemented free
subspaces are central to the work of Godefroy and Kalton
\cite[\S2]{GodefroyKalton03}.

A Banach space $X$ has the \emph{approximation property} (AP) if for every
norm-compact set $C\subseteq X$ and every $\varepsilon>0$ there is a bounded
finite-rank linear operator $T:X\to X$ with
$\sup_{u\in C}\norm{Tu-u}<\varepsilon$. For $\lambda\ge1$, it has
$\lambda$-BAP if these operators can always be chosen with
$\norm{T}\le\lambda$. The bounded approximation property (BAP) means
$\lambda$-BAP for some finite $\lambda$, and the metric approximation
property (MAP) means $1$-BAP. In AP the norm may depend on $C$ and
$\varepsilon$; in BAP one bound must work for all of them. A metric space
is \emph{proper} if its closed bounded balls are compact.

\begin{theorem}\label{thm:main}
There is a countable, unbounded, nonproper pointed metric space $(M,d,o)$
such that $d(x,y)\ge1$ for distinct points and $\F(M)$ has AP but fails
$\lambda$-BAP for every finite $\lambda\ge1$.
\end{theorem}

The distinction between discrete and uniformly discrete matters here.
Smith gives countable complete discrete spaces whose free spaces have AP
and fail BAP, without asserting a common positive lower bound on all
nonzero distances \cite[Theorem~1.3 and the surrounding discussion]{Smith26}.
\Cref{thm:main}
supplies that lower bound. The example is necessarily outside the countable
proper case: Dalet proved MAP for free spaces over countable proper metric
spaces \cite[Theorem~2.6]{Dalet15}. Our example contains an infinite separated
subset in a bounded ball.
Uniform discreteness also imposes strong sequential structure: C\'uth and
Kalenda prove a quantitative strengthening of the Schur property for these
free spaces \cite[Theorem~1.1]{CuthKalenda26}, so weakly convergent sequences
converge in norm. Their April 2026 introduction still records the BAP
question addressed here.

The same example has a nonlinear consequence. Kalton asked whether every
separable Banach space is \emph{approximable} \cite[Problem~1, p.~1260]{Kalton12}.
For such a space $X$, this means that there is an equi-uniformly continuous
sequence of maps $\psi_n:X\to X$, each with relatively compact range, such
that $\psi_n(x)\to x$ for every $x\in X$. Equi-uniform continuity means that
the choice of $\delta$ for a given $\varepsilon$ in the uniform-continuity
condition works for every $\psi_n$. The space $X=\F(M)$ in \Cref{thm:main}
is not approximable, although it has AP. \Cref{sec:nonlinear} explains the
short implication from Kalton's net theorem that gives this consequence.

\subsection{A finite obstruction in each bounded block}

The main construction concerns a prescribed operator norm bound $p$.
It produces a bounded metric space and a finite set of points, called
anchors, at which every finite-rank operator of norm at most $p$ makes a
definite error. The distinction from earlier finite extension obstructions
is important. For a finite pointed metric pair $A\subseteq N$ containing
the basepoint, a linear extension operator
$E:\Lip_0(A)\to\Lip_0(N)$ must satisfy $(Ef)|_A=f$ for every $f$.
Brudnyi and Brudnyi give finite subsets of integer $\ell_1$ cubes for which
the least possible norm of an exact linear extension operator grows without bound
\cite[preprint Lemma~10.5, pp.~46--47]{BrudnyiBrudnyi07}. Godefroy and Ozawa
\cite[Proposition~7]{GodefroyOzawa14} obtain arbitrarily large lower bounds for
finite metric pairs from a compact free space failing AP. By
finite-dimensional duality, an extension operator corresponds, with the
same norm, to a projection from $\F(N)$ onto the span of the evaluations
from $A$. These bounds therefore concern projections with a prescribed
range. On the finite-dimensional space $\F(N)$, the identity remains a
finite-rank operator of norm at most one. Our block is infinite, and its
fixed finite anchor set obstructs approximation by every operator of norm
at most $p$ whose range is finite-dimensional.

\begin{theorem}\label{thm:blocks}
For each integer $p\ge1$, there are a bounded countably infinite pointed
metric space $(M_p,d_p,o)$ and a finite set $A_p\subseteq M_p$ containing
$o$ such that
\[
 d_p(x,y)\ge1\quad(x\ne y),\qquad
 \operatorname{diam}M_p\le B_p:=300p^2+150p+5,
\]
and every bounded finite-rank linear operator
$T_0:\F(M_p)\to\F(M_p)$ with $\norm{T_0}\le p$ satisfies
\[
 \max_{a\in A_p}\norm{T_0\delta_a-\delta_a}\ge\frac12.
\]
\end{theorem}

Each individual block nevertheless has BAP. The coefficient identification
for bounded uniformly discrete spaces appears in
\cite[Theorem~4.14]{AACD20}, with the bounds used here in
\cite[Lemma~3.4]{CuthKalenda26}. Our proof in \Cref{prop:coefficients} identifies $\F(M_p)$ with
$\ell_1(M_p\setminus\{o\})$, and coordinate truncations give
$2B_p$-BAP alongside the obstruction at $p$. Transporting the free-space
norm to real $\ell_1$ gives the following quantitative consequence. It
also answers a separate question: does every equivalent norm on real
$\ell_1$ have MAP? Kalton observed that a positive answer would imply BAP
for every uniformly discrete free space \cite[p.~1260]{Kalton12}. The case
$p=1$ below gives a direct negative answer.

\begin{corollary}\label{cor:renorming}
For each integer $p\ge1$, real $\ell_1$ has an equivalent norm
$\norm{\cdot}_p$ such that
\[
 \tfrac12\norm{c}_1\le\norm{c}_p\le B_p\norm{c}_1,
\]
and $(\ell_1,\norm{\cdot}_p)$ has $2B_p$-BAP but fails $p$-BAP.
For $p=1$, this norm fails MAP.
\end{corollary}

To pass from the blocks to \Cref{thm:main}, identify their
basepoints and put
\[
 d(x,y)=d_p(x,o)+d_q(y,o)
 \quad\text{if }x\in M_p\setminus\{o\},\ y\in M_q\setminus\{o\},\ p\ne q.
\]
Each block then has a pointed $1$-Lipschitz retraction from the resulting
metric wedge. On free spaces, inclusion and retraction induce contractions
whose composite is the identity. Thus a BAP bound for the wedge would be
a bound for every block, contradicting \Cref{thm:blocks}.

The same decomposition proves AP directly. First retain finitely many
blocks by a contractive projection that approximates the chosen compact
set. Then approximate its images in those blocks, using their individual
BAPs. Only finitely many operator norms occur at this stage, so the
resulting operator is bounded. The bound is allowed to grow as more
blocks are retained.

\subsection{The obstruction: signed measures and opposing tests}

The block itself consists of two kinds of points: vertices of finite source
graphs and anchors that join those graphs into a bounded metric space.
After a finite-rank operator is given, we choose a finite subset $S$ of
these same points to support the range of a nearby operator. Thus $S$ may
contain source vertices or anchors and may depend on the operator.

The source graphs have a fixed large degree. For every fixed radius, the
frequencies of their colored neighborhoods approach the corresponding
probabilities in a regular tree of that degree with independent uniform
vertex colors. A vertex is labeled by its
capped distances to the color classes. These labels vary by at most one
across an edge, and each label has only a bounded number of coordinates
below the cap. Vertices with the same label attach to the same anchor
by a long edge. There are only finitely many possible labels, hence
finitely many anchors. Joining the anchors makes the whole block bounded,
while sufficiently small metric balls remain inside a single source
graph and contain a uniformly bounded number of points.

Suppose a finite-rank operator of small norm approximates the
anchor evaluations. An arbitrarily small perturbation gives an operator
$T$ with $\norm{T}\le2p$ whose range has a common finite support
$S\subset M_p$. Adding a coefficient at the
basepoint allows us to write its values as signed measures $\mu_x$ on
$S$, each of mass one and uniformly bounded total variation. For the
height $L=10p$, a nonnegative bump $b_x:M_p\to[0,L]$ detects support points
whose labels are close to the label of $x$. It takes its maximum $L$ at the
anchor $a(x)$ attached to $x$ and has a small Lipschitz constant, in
particular at most one. The operator norm controls the comparison from
$x$ to $a(x)$, while the anchor approximation controls the value there.
The small slope of $b_x$ makes their combined loss from the anchor value
at most $L/4$. The affine representation therefore gives
\[
 (T\delta_x)(b_x-b_x(o))+b_x(o)=\int b_x\,d\mu_x\ge\frac{3L}{4}.
\]
No positivity of $\mu_x$ is assumed.

A source walk is \emph{nonbacktracking} if it never immediately reverses
an edge incidence. The color structure makes it rare for a source vertex
to start a long nonbacktracking walk that can be tracked by nearby support
points with compatible labels. Such a walk is \emph{matched} if one can
choose a support point at each step with a compatible label, while
consecutive support points stay within distance $2L$. For a prescribed
support point, compatibility allows only a
small set of source colors. A bound on small metric balls limits the
possible next support points. On the limiting tree, the source colors
along a nonbacktracking walk are independent. Choosing many colors
therefore makes long matched walks rare.

Call a vertex \emph{good} if no matched walk of the chosen long length
starts there, and \emph{bad} otherwise. At a good vertex, record the longest matched walk
beginning in each outgoing direction, for each fixed
initial support point. Declare a direction exceptional only when its
record is strictly larger than every other record. There is at most
one such direction per support point. If nearby support points match
the two ends of an edge whose endpoints are good, at least one of the two directions along that
edge must be exceptional: otherwise prepending the edge to a longest
walk in another direction would force each of two finite records to be
larger than the other.

For an oriented incidence $e:x\to y$, let $E_e$ be the support points for
which $e$ is exceptional at $x$; at a bad vertex we put every support point
in $E_e$. We prescribe $b_x$ on $S\setminus E_e$ and $-b_y$ on
$S\setminus E_{\bar e}$, where $\bar e$ is the reverse incidence. Within
each set the prescribed values satisfy the Lipschitz inequality because
the bumps do. For a cross pair $j\in S\setminus E_e$ and
$j'\in S\setminus E_{\bar e}$, the record argument shows that if
$d_p(j,j')<2L$, then both $b_x(j)$ and $b_y(j')$ vanish. If the distance
is at least $2L$, their sum is bounded by that distance. Thus all the
prescribed values together are $1$-Lipschitz. Scalar extension followed by
clipping gives a $1$-Lipschitz test $f_e:M_p\to[-L,L]$. Its pairing satisfies
\[
 \int f_e\,d(\mu_x-\mu_y)
 \ge \frac{3L}{2}
 -2L\bigl(|\mu_x|(E_e)+|\mu_y|(E_{\bar e})\bigr).
\]
Here $|\mu_x|(E_e)$ is the sum of the absolute values of the coefficients
of $\mu_x$ on $E_e$. Each support point belongs to at most one exceptional
set at a good vertex. Averaging over all oriented incidences therefore
divides the corresponding variation loss by the graph degree; the small
fraction of bad vertices controls the remaining loss. Since source-edge
endpoints have $d_p$-distance at most one, the affine pairing bounds each
test pairing in absolute value by $\norm{T}\le2p$. The averaged lower bound
contradicts this upper bound.

The order of choices is essential. The entire deterministic graph
sequence, metric block, and anchors are fixed before the operator and
its support $S$. We then choose the walk length and one sufficiently
large source graph. Local convergence applies to every finite colored
event, including the event specified by this already fixed support.
There is no assumption that the operator was chosen independently of
the graph sequence.

\subsection{Methods and organization}

For the graph construction we use the independent-permutation model
described by Friedman \cite{Friedman04} and prove the local convergence
we need directly. The empirical rooted-neighborhood viewpoint is that
of Benjamini and Schramm \cite{BenjaminiSchramm01}. The scalar extension
step is an instance of McShane's classical extension principle
\cite{McShane34}; we give the needed formula for finite data and its proof.
A related comparison of nonbacktracking continuation lengths is used to
select parent directions in the graph algorithm of Ben-David, Childs,
Gily\'en, Kretschmer, Podder, and Wang
\cite[\S6.1.2]{BenDavidEtAl24}. Our continuation records also depend on
the support state and are used to construct the scalar tests.

\Cref{sec:reduction} proves the coefficient estimate, the renorming
corollary, and the wedge reduction, assuming \Cref{thm:blocks}, and derives
the nonlinear consequence from Kalton's theorem.
\Cref{sec:blocks} constructs the colored graphs and bounded block.
\Cref{sec:operators,sec:walks,sec:tests} prove \Cref{thm:blocks}:
they pass from an arbitrary finite-rank operator to signed measures,
bound matched walks, and derive the contradiction by scalar tests and
incidence averaging.
